\section{Discussion and remaining questions}\label{sec:discussion}
The common design principle is to optimize the prefix curve through a
weighted sum of expected squared discrepancies. For lattice sequences,
dyadic difference multiplicities make this objective tractable. The Newo
invariance result shows that changing the common scalar digit ordering
within the stated class cannot improve this particular criterion for a
fixed generator: every prefix value is already unchanged. Such scrambling
can still change the appearance of a prefix and its integration error for
a particular function. These are separate effects from WSESD optimization.

For Kronecker sequences, the natural decomposition is by integer lags.
It leads to a comparable product-kernel CBC objective, but neither the
lattice search's acceleration nor its available theory transfers merely
because both methods are called CBC. The newer Kronecker construction
should be specified on its own terms, including its candidate family,
search procedure, and finite-horizon or asymptotic guarantees.

The conference abstract proposed arbitrary sample sizes and missing data
as motivations. In this paper, the main results concern all prefixes up to
$N$, whereas~\eqref{eq:removal} gives a limited deletion-stability statement.
A construction specifically optimized over missing-point patterns would
need an additional criterion. Likewise, this draft does not extend the
construction horizon without rebuilding the generator or establish
a stopping rule based on observed integrand values.

\subsection*{Questions for the discussion with Anders}
\codex{[Review agenda for the authors:]}
\codex{\begin{enumerate}[leftmargin=*,itemsep=3pt]
\item Confirm the finite-precision $\pi$ definition and proof in
Section~\ref{sec:newo}, and whether a single infinite-precision family is
needed for the intended implementation.
\item Supply the distinct, newer Kronecker CBC: exact candidate parameters,
objective or surrogate, update rule, cost, and available bounds. Distinguish
it from the \texttt{ANDERS\_CBC} generator used in the talk.
\item Identify a stable implementation and a reproducible figure-generation
path for the lattice and Kronecker searches, including the random-shift
convention in the discrepancy evaluator.
\item Agree on a small numerical study that fits the 16-page limit:
all-prefix comparisons, preferred-size tradeoffs, and a representative
integration example if space permits.
\item Confirm attribution of the new results and the eventual
proceedings format before submission.
\end{enumerate}
These questions identify work remaining in the draft rather than claims
that the numerical results or the newer algorithm have already been verified.}
